\clearpage\begingroup\NoSourceHead\SourceLeaf{12}
\begin{center}
[II.]\\[.7em]
BERICHT ÜBER DIE ART, WIE DIE HANNOVERSCHEN NORMAL-\\
PFUNDE DARGESTELLT SIND.\\[1em]
\textbf{Allgemeine Vorerinnerung.}
\end{center}
\noindent\hspace*{\parindent}Zwei Gewichtsstücke sind nur insofern als gleich zu betrachten, als sie
gleiche Quantitäten ponderabler Masse enthalten; wäre es thunlich, Abwägungen
im leeren Raume auszuführen, so würden dann zwei gleiche Gewichtsstücke
auf der Waage einander das Gleichgewicht halten. Bei der
Wägung in der Luft, wo bekanntlich jeder Körper soviel am Gewicht verliert,
als er Luft verdrängt, würde dieser Gewichtsverlust auf das Gleichgewicht
keinen Einfluss haben, wenn jedes der beiden Gewichtsstücke gleichviel verlöre.
Allein da dieser Gewichtsverlust von dem Raumvolumen, welches jeder
Körper einnimmt, mithin bei gleichem Masseninhalt von der specifischen
Schwere abhängt, so \textls[120]{erscheinen} wirklich gleiche aber an specifischer Schwere
ungleiche Körper auf der Waage als ungleich; und umgekehrt, aus dem Gleichgewicht
auf der Waage folgt noch nicht die Gleichheit der Gewichte, oder
vielmehr, es folgt daraus die Ungleichheit der Gewichte, wenn man weiss,
dass sie von verschiedener specifischer Schwere sind. Es erhellt daraus, dass
alle Wägungsvergleichungen, wobei eine grosse Schärfe verlangt wird, erst
einer \textls[120]{Reduction} auf den leeren Raum bedürfen, wozu die Kenntniss der
specifischen Schweren der gewogenen Körper das erste Erforderniss ist; auch
muss, wo es irgend genau genommen werden soll, der jedesmalige Dichtigkeitszustand
der Luft, nach Barometer- und Thermometer-Stand gehörig berücksichtigt
werden.

Für die Bedürfnisse des gewöhnlichen Verkehrs sind zwar diese Unterschiede
\clearpage\SourceLeaf{13}\HistoricalHead{8}{L}{PHYSIK, AMTLICHE BERICHTE.}
\noindent als ganz unerheblich zu betrachten: allein für feinere wissenschaftliche
Zwecke, oder für Darstellung von Normal-Gewichten darf von einer Ignorirung
dieser Unterschiede nicht die Rede sein: sie müssen, so scharf wie es die zu
Gebote stehenden Hülfsmittel verstatten, berücksichtigt werden. Zwei an
Masseninhalt ganz gleiche Pfunde, eines aus Platin, das andere aus Messing,
müssen auf der Waage in der Luft als ungleich erscheinen, das erstere um
40—45 Milligramme schwerer als das letztere: umgekehrt, ein Platinpfund,
welches einem Messingpfunde auf der Waage genau das Gleichgewicht hält,
ist der Wahrheit nach um 40—45 Milligramme leichter; ja, es ist klar, dass
dasselbe Platinpfund, welches in Luft von einer bestimmten Dichtigkeit dem
Messingpfunde gleich erschien, diesem ungleich erscheinen wird, wenn die
Wägung in Luft von einer anderen Dichtigkeit geschieht. Die Waage für
sich allein, ohne Reduction auf den leeren Raum, gibt also, wenn die gewogenen
Körper von ungleicher specifischer Schwere sind, ungleiche Resultate,
andere an einem höhern, andere an einem tiefer liegenden Orte u.s.w., und
jene Reduction auf den leeren Raum wird daher allgemein, wo auf Schärfe
Anspruch gemacht wird, wie eine sich von selbst verstehende Bedingung betrachtet.
Dies gilt aber natürlich nicht bloss, wenn von Gewichtsstücken aus
ungleichem Metall, sondern auch, wenngleich in geringerem Maasse, wenn von
solchen die Rede ist, die aus einerlei oder vielmehr aus gleichnamigem Metall
bestehen. Messing z. B. variirt in der specifischen Schwere so sehr, dass zwei
Messingpfunde von gleichem Gewichtsinhalt auf der Waage um 5 Milligramme
verschieden erscheinen können. Bei Messing ist dies desto weniger auffallend,
da es eine Mischung aus zwei Metallen ist; ähnliches zeigt aber die Erfahrung
auch bei Gewichtsstücken aus Platin, vermuthlich nur in Folge der ungleichen
chemischen Reinheit.

Wenn demnach zu einer genauen Copirung eines gegebenen Gewichtsstückes
nach seinem wahren Masseninhalt, oder zu einer genauen Vergleichung
zweier Gewichtsstücke untereinander, die Reduction der gemachten Wägungen
auf den leeren Raum durch Calcül nothwendig ist, so ist klar, dass dazu
eine genaue Kenntniss der specifischen Schwere unumgänglich erfordert wird;
zur Gewinnung dieser Kenntniss ist das Abwägen in destillirtem Wasser das
einzige Mittel, welches nur unter selten vorkommenden Bedingungen durch
ein anderes Verfahren theilweise ersetzt werden kann.

\clearpage\SourceLeaf{14}\HistoricalHead{9}{R}{ÜBER DIE HANNOVERSCHEN NORMALPFUNDE.}
Es darf nun hier nicht verschwiegen werden, dass in dieser Beziehung
sowohl der französischen Commission, die das französische, als der preussischen,
die das preussische Gewichts-System regulirt hat, der Vorwurf zur Last fällt,
dass sie versäumt haben, von denjenigen Gewichtsstücken, die als die gesetzlichen
Normal-Einheiten deponirt sind, vorher die specifische Schwere direct
und scharf zu bestimmen. Hinterher ist dies nicht mehr zulässig wegen der
Besorgniss, dass diese Urgewichte durch Abwägen im Wasser eine wenn auch
geringe dauernde Veränderung erleiden könnten: man hätte dieses \textls[120]{vorher}
thun sollen, ehe sie fertig regulirt und ihnen ein gesetzlicher Charakter beigelegt
war. Als eine Entschuldigung mag angesehen werden, dass man in
jener Zeit, wo das französische Maass- und Gewicht-System regulirt wurde,
noch nicht ganz so hohe Anforderungen stellte, wie gegenwärtig: dass die
Waagen, die man in Paris und Berlin gebrauchte, das einzelne Milligramm
schwerlich mit \textls[120]{Zuverlässigkeit} ergaben, und dass man sich dabei beruhigte,
Platin sei Platin, und Messing Messing. Es entstehen aber aus jenem Umstande
gegenwärtig eigenthümliche Schwierigkeiten für denjenigen, der ganz
dem heutigen Zustande der Hülfsmittel Genüge leisten will, Schwierigkeiten,
die mir selbst erst im Fortgange des Geschäfts ganz fühlbar geworden sind,
während es mir gelang, der dabei gebrauchten an sich schon vortrefflichen
\textsc{Repsold}schen Waage durch eigenthümliche Verbesserungen und eigenthümliche
Behandlungsart eine Vollkommenheit zu geben, wobei selbst kleine
Bruchtheile des Milligramms noch mit Zuverlässigkeit festgestellt werden.

In dem nachfolgenden Bericht über die Darstellung der Hannoverschen
Normalpfunde habe ich zur Abkürzung, und um Verwirrung zu vermeiden,
die verschiedenen darin erwähnten Gewichtsstücke mit Buchstaben bezeichnet,
die ich hier in Einer Uebersicht zusammenstelle.
\begin{list}{}{\setlength{\leftmargin}{2.4em}\setlength{\labelwidth}{1.6em}\setlength{\labelsep}{.5em}\setlength{\itemsep}{0pt}\setlength{\parsep}{0pt}\setlength{\topsep}{0pt}}
\item[A.] das französische gesetzliche Normalkilogramm aus Platin, welches seit
22. Junius 1799 im Reichsarchiv unter vier Schlössern aufbewahrt
und nur in ganz ausserordentlichen Fällen zugänglich gemacht wird.
Hiervon stammen alle andern Kilogramme als Copien, oder Copien
von Copien, oder in noch mehreren gar nicht mehr nachzuweisenden
Zwischenstufen, ab; mittelbarer Weise auch das preussische Pfund.
\item[B.] ein Kilogramm aus Platin von \textsc{Gambey} verfertigt, im Besitz des Conferenzraths
\textsc{Schumacher} in Altona. Diess ist, wie weiter unten erwähnt
\HistoricalFoot{XI.}{2}
\clearpage\SourceLeaf{15}\HistoricalHead{10}{L}{PHYSIK, AMTLICHE BERICHTE.}
\noindent werden wird, die zuverlässigste Copie, die in der ganzen Welt
existirt.
\item[C.] ein Kilogramm aus Messing (in zwei Halbkilogrammenstücken), welches
ich durch \textsc{Repsold} in Hamburg nach B habe machen lassen.
\item[D.] das preussische Urpfund, welches, wenn ich nicht irre, bei dem Ministerium
des Handels und der Finanzen aufbewahrt wird.
\item[E.] eine Copie von D, welche ich 1836 durch den Oberbergrath \textsc{Schaffrinsky}
in Berlin habe machen lassen.
\item[F.] ein preussisch-hannoversches Pfund von \textsc{Repsold} im Besitz des Conferenzrathes
\textsc{Schumacher}.
\item[G.] und H. die beiden von \textsc{Meyerstein} angefertigten Normalpfunde, und
zwar H dasjenige, welches jetzt nach Hannover abgesandt wird, G das
andere, welches in Göttingen verbleiben und früherer Bestimmung
zufolge bei der Societät der Wissenschaften hinterlegt werden soll.
\end{list}
\HistoricalHead{10}{L}{PHYSIK, AMTLICHE BERICHTE.}
\noindent\hspace*{\parindent}Die Gewichte D—H sind alle von Messing und vergoldet.

Mit denselben Buchstaben sind diese Gewichte in den Kostenrechnungen,
soweit sie darin vorkommen, bezeichnet.
\begin{center}\rule{1.35in}{.35pt}\end{center}
\noindent\hspace*{\parindent}Das preussische Urpfund D ist nach dem ausführlichen von \textsc{Eytelwein}
in den Abhandlungen der Berliner Akademie der Wissenschaften für 1825 gegebenen
Berichte[\textsuperscript{*})]\begingroup\renewcommand{\thefootnote}{}\footnotetext{[\textsuperscript{*}) \textsc{Eytelwein}, \textit{Über die Prüfung der Normal-Maasse und Gewichte für den königlich-preussischen Staat und ihre Vergleichung mit den französischen Maassen und Gewichten}, Abhandlungen der mathem. Klasse der K. Akad. d. Wissensch. zu Berlin (1825) 1828, S. 1.]}\endgroup\ so angefertigt, dass man, so genau es mit den vorhandenen
Hülfsmitteln möglich war, 467711 Milligramme, d. i. Milliontheile der französischen
Gewichtseinheit darzustellen sich bestrebt hat. Es waren dazu zwei
Kilogramme, eines aus Platin, das andere aus Messing, benutzt, zwischen
denen sich (nach Reduction der Wägungen auf den leeren Raum) eine Differenz
von einem Milligramm fand; der Berichterstatter erklärt selbst, dass die
in Paris gemachten Vergleichungen zwischen jenen beiden Copien und dem
Original ein einzelnes Milligramm verbürgende Schärfe nicht gehabt haben
möchten. Jener Werth von 467711 Milligrammen für das ganze, oder von
\(238855\frac{1}{2}\) Milligrammen für das halbe preussische Pfund (die cölnische Münzmark)
ist als officiell zu betrachten.
\par\clearpage\endgroup
